\documentclass[10pt,a4paper]{amsart}
\usepackage[margin=19mm]{geometry}
\usepackage{amsmath,amssymb,mathtools}
\usepackage[scr=rsfs,cal=euler]{mathalfa}
\usepackage{xfrac}
\usepackage[unicode,hidelinks,bookmarks=false]{hyperref}
\newcommand{\ie}{i.e.\ }
\newcommand{\cf}{cf.\ }
\newcommand{\resp}{respectively\ }
\newcommand{\Subsection}{\subsection}
\providecommand{\nobreakdash}{\nobreak\hbox{-}}
\setlength{\emergencystretch}{2em}
\multlinegap0pt
\usepackage{bm}
\usepackage{bbm}
\usepackage{dsfont}
\usepackage{xspace}
\usepackage{cancel}
\usepackage{mathtools}
\usepackage{cleveref}
\usepackage{amsthm}
\makeatletter
\@ifundefined{lemma}{\newtheorem{lemma}{Lemma}}{}
\@ifundefined{theorem}{\newtheorem{theorem}{Theorem}}{}
\makeatother
\newcommand{\E}{E}
\newcommand{\R}{\mathbb{R}}
\newcommand{\X}{\mathbb{X}}
\newcommand{\calC}{\mathcal{C}}
\newcommand{\calT}{\mathcal{T}}
\newcommand{\dd}{\mathrm d}
\newcommand{\deq}{\stackrel{\mbox{\scriptsize d}}{=}}
\newcommand{\prob}{{\rm pr}}
\renewcommand{\mid}{\ensuremath{\,|\,}}
\title[Palm distributions of superposed point processes for statist]{Palm distributions of superposed point processes for statistical inference}
\author{Mario Beraha, Federico Camerlenghi, Lorenzo Ghilotti}
\date{Source: arXiv:2508.20924v2 (CC BY 4.0)}
\begin{document}
\maketitle

\noindent\emph{Source and licence.} Palm distributions of superposed point processes for statistical inference. Version arXiv:2508.20924v2 (CC BY 4.0), distributed under the Creative Commons Attribution 4.0 International licence, as recorded in the arXiv licence metadata for this exact version and shown on the abstract page https://arxiv.org/abs/2508.20924v2. Full rights notice: licenses/arxiv-2508.20924-CC-BY-4.0.txt.\par\smallskip

\noindent\textit{Editorial scope.} Counted unit: the Theorem whose source label is \texttt{None} in the article (source0.tex lines 1442--1459, source label \texttt{None}), delivered together with the article's own complete original proof of it (source0.tex lines 1460--1515, one \texttt{proof} environment of 56 lines). No mathematics is written, repaired or re-derived: statement and proof are the article's own text line for line. Required context retained uncounted: thm:red\_palm\_sncp (lines 384--399); thm:janossy\_general (lines 572--581); lemma:k\_factorial\_sncp (lines 1320--1326); prop:prior\_partition\_sncp (lines 1335--1341). Every definition, display and label of those lines is retained unchanged. Omitted from this package and not redistributed: 1--383, 400--571, 582--1319, 1327--1334, 1342--1441, 1516--1686. Nothing inside a retained excerpt is omitted or altered: every retained body is delivered exactly as the article's lines are written, so no omission receipt of any kind accompanies this package. Citations: the retained excerpts carry no citation command at all, so no bibliography entry of the article is transcribed and no reference is redistributed. Reference adaptation: none was needed and none was made; every reference inside the delivered unit resolves inside it, so no unresolved reference or citation is produced. Printed slips kept literal: no printed word, symbol, hypothesis or number of the counted statement or of its proof was altered. Layout and class: the article's own class is \texttt{article} with packages including geometry, babel, inputenc, fontenc, amsthm, amsmath, amssymb, mathrsfs, dsfont, mathtools, booktabs, subcaption, bbm, multirow; the wrapper is the lane's pinned \texttt{amsart} editorial wrapper at 10pt on A4, which restates the theorem-like environments the retained text needs and adds the article's own macro definitions that the retained text needs (\texttt{\textbackslash E}, \texttt{\textbackslash R}, \texttt{\textbackslash X}, \texttt{\textbackslash calC}, \texttt{\textbackslash calT}, \texttt{\textbackslash dd}, \texttt{\textbackslash deq}, \texttt{\textbackslash mid}, \texttt{\textbackslash prob}), with extra wrapper packages (bm, bbm, dsfont, xspace, cancel, mathtools, cleveref). The delivered numbering is the unit's own. Assets: no retained span contains a figure environment, a graphics-inclusion command, a PDF-page inclusion, a table or a TikZ picture; no image file is copied into this package, the package declares no asset dependency and no image file is redistributed with it. Licence: version arXiv:2508.20924v2 of this article is distributed under the Creative Commons Attribution 4.0 International licence, as recorded in the arXiv licence metadata for this exact version and shown on the abstract page https://arxiv.org/abs/2508.20924v2. A full rights notice is delivered at licenses/arxiv-2508.20924-CC-BY-4.0.txt. Characters: every retained excerpt is pure ASCII, so the delivered engine, XeLaTeX with its default Latin Modern text font, reports no missing character.
\begin{theorem}\label{thm:red_palm_sncp} 
Let $\Phi \sim \textsc{sncp}(\kappa, \nu)$ and $\underline{x} = (x_1, \ldots, x_k) \in \X^k$. Then, the reduced Palm version of $\Phi$ at $\underline{x}$ admits the following representation
\begin{equation}\label{eq:thesis_palm_sncp}
    \Phi^!_{\underline{x}}\mid \underline{T} \deq \Phi + \sum_{h = 1}^{|\underline{T}|} \Phi_{\zeta_{{\underline{x}}_h}},\qquad \prob(\underline{T} = \underline{t}) \propto \prod_{h=1}^{|\underline{t}|} \eta(\underline{x}_h),
\end{equation}
where $\underline{T} := (T_1,\ldots,T_k)$ are latent indicators describing a partition of $\underline{x}$ into $|\underline{T}|$ clusters and
\begin{equation*}
    \begin{aligned}
        \Phi_{\zeta_{{\underline{x}}_h}} \mid \zeta_{{\underline{x}}_h} = (\theta_{{\underline{x}}_h}, \gamma_{{\underline{x}}_h}) &\sim \textsc{pp}(\gamma_{{\underline{x}}_h} \kappa(x; \theta_{{\underline{x}}_h}) \,\dd x)\\
        \zeta_{{\underline{x}}_h}  &\sim f_{{\underline{x}}_h}(\dd \theta\, \dd \gamma) \propto \gamma^{n_h} \prod_{x_j \in 
        \underline{x}_h}\kappa(x_j;\theta) \nu(\dd \theta\, \dd \gamma),
    \end{aligned}
\end{equation*}
where $\underline{x}_h = (x_j: T_j = h)$ and $n_h$ is the cardinality of $\underline{x}_h$.
Finally, the processes $\Phi$ and $\Phi_{\zeta_{{\underline{x}}_h}}$, $h = 1,\ldots, |\underline{T}|$, are mutually independent conditionally to $\underline{T}$.
\end{theorem}

\begin{theorem}[Janossy measures of finite point processes]\label{thm:janossy_general}
Let $\Phi$ be a finite point process on $\X$. Then, its Janossy measures satisfy
\[
J_k(\dd x_1\,\ldots\, \dd x_k) = \prob(\Phi^!_{(x_1,\ldots,x_k)}(\X\setminus \{x_1,\ldots,x_k\} ) = 0)  M^{(k)}_{\Phi}(\dd x_1\,\ldots\, \dd x_k).
\]
Moreover, if $\Phi$ is a simple point process, it simplifies as
\[
J_k(\dd x_1\,\ldots\, \dd x_k) = \prob(\Phi^!_{(x_1,\ldots,x_k)}(\X) = 0)  M^{(k)}_{\Phi}(\dd x_1\,\ldots\, \dd x_k).
\]
\end{theorem}

\begin{lemma}\label{lemma:k_factorial_sncp}
Let $\Phi \sim \textsc{sncp}(\kappa, \nu)$. The $k$-th order factorial moment measure, denoted with $M^{(k)}_{\Phi}$, is equal to 
\[
M^{(k)}_{\Phi}(\dd \underline{x}) = \sum_{\underline{t} \in \calT_k} \prod_{h=1}^{|\underline{t}|} \eta(\underline{x}_h)\dd \underline{x},
\]
where $\underline{x}_h = (x_j: t_j = h)$, for $h=1,\ldots,|\underline{t}|$.
\end{lemma}

\begin{theorem}\label{prop:prior_partition_sncp}
Consider $\Phi \sim \textsc{sncp}(\kappa, \nu)$, for some $\kappa$ and $\int_{\X\times \R_+} (1-e^{-\gamma} ) \nu(\dd\theta \, \dd \gamma) < \infty$. The joint distribution on the number of points and the latent partition induced by $\Phi$ is
\[
\prob(N, \{\calC_1,\ldots, \calC_C\}) = \frac{1}{N!} e^{- \int_{\X\times\R_+}(1-e^{-\gamma}) \nu(\dd\theta \, \dd \gamma)} \prod_{h = 1}^C \int_{\X\times \R_+} e^{-\gamma} \gamma^{n_h} \nu(\dd\theta \, \dd \gamma), 
\]
where $n_h = | \calC_h |$, $N = \sum_{h=1}^C n_h$. Clearly, this law does not depend on the kernel $\kappa$.
\end{theorem}

\begin{theorem}
    Let $\Phi \sim \textsc{sncp}(\kappa, \nu)$ such that $\Phi(\X) < \infty$ almost surely. Then, its Janossy density equals
    \begin{equation*}
        \begin{aligned}
            j_k(x_1,\ldots,x_k) = \exp\left\{ - \int_{\X\times \R_+} ( 1 - e^{-\gamma}) \nu(\dd \theta\, \dd \gamma) \right\} \sum_{\underline{t} \in \calT_k} \prod_{h=1}^{|\underline{t}|} \int_{\X\times \R_+} e^{-\gamma} \gamma^{n_h} \prod_{j: t_j = h} \kappa(x_j ; \theta) \nu(\dd \theta \,\dd \gamma).
        \end{aligned}
    \end{equation*}
     
 
 Moreover, if $\nu(\dd \theta\, \dd \gamma) = \rho(\dd \gamma) G_0(\dd \theta) $, where $G_0$ is a probability measure on $\X$, then 
\begin{equation*}
    j_k(x_1,\ldots,x_k) = k! \, \prob( \Phi(\X) = k) \,\E\left[\prod_{h=1}^{|\underline{T}|} \int_{\X}\prod_{j: T_j = h} \kappa(x_j ; \theta) G_0(\dd \theta)\right],
\end{equation*}
where the expectation is taken with respect to the indicators $\underline{T} = (T_1, \ldots, T_k) \in \calT_k$ with distribution
\begin{equation*}
\prob(\underline{T} = \underline{t}) \propto  \prod_{h=1}^{|\underline{t}|} \int_{\R_+} e^{-\gamma} \gamma^{n_h} \rho(\dd \gamma) .    
\end{equation*}
\end{theorem}

\begin{proof}
The proof follows from specializing \Cref{thm:janossy_general} for a \textsc{sncp}. In particular, \Cref{thm:janossy_general} states that
\begin{equation}\label{eq:janossy_two_terms}
    J_k(\dd x_1\,\ldots\, \dd x_k) = \prob(\Phi^!_{(x_1,\ldots,x_k)}(\X) = 0)  M^{(k)}_{\Phi}(\dd x_1\,\ldots\, \dd x_k).
\end{equation}

Focusing on the first term of \eqref{eq:janossy_two_terms}, and denoting with $\underline{x} = (x_1,\ldots,x_k)$,
\begin{equation}\label{eq:ev_janossy_start}
        \prob \left( \Phi^!_{\underline{x}}(\X) = 0 \right) = \E \left[ \prob \left( \Phi^!_{\underline{x}}(\X) = 0 \mid \underline{T} \right) \right] = \E \left[ \prob \left( \Phi(\X) = 0 \right) \prod_{h=1}^{|\underline{T}|} \prob \left( \Phi_{\zeta_{\underline{x}_h}}(\X) = 0 \right)  \right], 
\end{equation}
where $\underline{T}$ and the $\Phi_{\zeta_{\underline{x}_h}}$'s are introduced in \Cref{thm:red_palm_sncp}. It follows that \eqref{eq:ev_janossy_start} writes as
\begin{equation*}
   \begin{aligned}
       \prob \left( \Phi^!_{\underline{x}}(\X) = 0 \right) &= \E\left[ \prob \left( \Phi(\X) = 0 \mid \Lambda \right)\right] \E \left[ \prod_{h=1}^{|\underline{T}|} \E\left[ \prob \left( \Phi_{\zeta_{\underline{x}_h}}(\X) = 0 \mid \zeta_{\underline{x}_h}\right)\right]  \right]\\
       &= \E\left[ e^{- \int_{\X} \int_{\X\times\R_+} \gamma \kappa(x; \theta) \Lambda(\dd \theta\, \dd \gamma) \dd x} \right] \E\left[ \prod_{h=1}^{|\underline{T}|} \E\left[ e^{- \int_{\X} \gamma_{\underline{x}_h} \kappa(x; \theta_{\underline{x}_h}) \dd x} \right] \right]\\
       &= \exp\left\{ - \int_{\X\times\R_+} \left(1 - e^{- \gamma } \right) \nu(\dd \theta \, \dd \gamma) \right\} \E\left[ \prod_{h=1}^{|\underline{T}|} \int_{\X\times\R_+} e^{-\gamma } f_{\underline{x}_h}(\dd\theta \, \dd \gamma) \right]\\
       &= \exp\left\{ - \int_{\X\times\R_+} \left(1 - e^{- \gamma  } \right) \nu(\dd \theta\, \dd \gamma) \right\} \sum_{\underline{t} \in \calT_k} \left[ \prod_{h=1}^{|\underline{t}|} \int_{\X\times\R_+} e^{-\gamma } f_{\underline{x}_h}(\dd\theta \, \dd \gamma) \right] \prob(\underline{T} = \underline{t}),
   \end{aligned} 
\end{equation*}
where $f_{\underline{x}_h}$ is defined in \Cref{thm:red_palm_sncp}.  Finally, plugging this last expression in \eqref{eq:janossy_two_terms} and using \Cref{lemma:k_factorial_sncp} for $M^{(k)}_{\Phi}$, we obtain
\begin{equation*}
    \begin{aligned}
        J_k(\dd x_1\,\ldots\,\dd x_k) &= \exp\left\{ - \int_{\X\times\R_+} \left(1 - e^{- \gamma } \right) \nu(\dd \theta \,\dd \gamma) \right\} \sum_{\underline{t} \in \calT_k} \prod_{h=1}^{|\underline{t}|} \eta(\underline{x}_h) \int_{\X\times\R_+} e^{-\gamma } f_{\underline{x}_h}(\dd\theta \, \dd \gamma)\; \dd x_1\,\ldots\,\dd x_k\\
        &= \exp\left\{ - \int_{\X\times\R_+} \left(1 - e^{- \gamma} \right) \nu(\dd \theta\, \dd \gamma) \right\} \sum_{\underline{t} \in \calT_k} \prod_{h=1}^{|\underline{t}|}  \int_{\X\times\R_+} e^{-\gamma } \gamma^{n_h} \prod_{j: t_j = h} \kappa(x_j; \theta) \nu(\dd\theta \, \dd \gamma) \; \dd x_1\,\ldots\,\dd x_k,
    \end{aligned}
\end{equation*}
and the thesis is proved.\\

If $\nu(\dd \theta\, \dd \gamma) = \rho(\dd \gamma) G_0(\dd \theta)$, where $G_0$ is a probability measure, then the Janossy density boils down to
\begin{equation}\label{eq:janossy_result_specialized}
\begin{split}
 & j_k(x_1,\ldots,x_k) \\
 \quad & = \exp\left\{ - \int_{\R_+} ( 1 - e^{-\gamma}) \rho(\dd \gamma) \right\} \sum_{\underline{t} \in \calT_k} \prod_{h=1}^{|\underline{t}|} \int_{\X} \prod_{j: t_j = h} \kappa(x_j ; \theta) G_0(\dd \theta) \int_{\R_+} e^{-\gamma} \gamma^{n_h} \rho(\dd \gamma).  
\end{split}
\end{equation}
Specializing the result in \Cref{prop:prior_partition_sncp} for $\nu(\dd \theta\, \dd \gamma) = \rho(\dd \gamma) G_0(\dd \theta) $, we obtain
\[
\prob(N, \{\calC_1,\ldots, \calC_C\}) = \frac{1}{N!} e^{- \int_{\R_+}(1-e^{-\gamma}) \rho(\dd \gamma)} \prod_{h = 1}^C \int_{\R_+} e^{-\gamma} \gamma^{n_h} \rho(\dd \gamma), 
\]
and the following holds true:
\begin{equation*}
\begin{aligned}
    \prob(N = k) &= \sum_{\{\calC_1,\ldots,\calC_C\}} \prob(N = k, \{\calC_1,\ldots, \calC_C\})\\
    &= \frac{1}{k!} e^{- \int_{\R_+}(1-e^{-\gamma}) \rho(\dd \gamma)} \sum_{\{\calC_1,\ldots,\calC_C\}} \prod_{h = 1}^C \int_{\R_+} e^{-\gamma} \gamma^{n_h} \rho(\dd \gamma)\\
    &= \frac{1}{k!} e^{- \int_{\R_+}(1-e^{-\gamma}) \rho(\dd \gamma)} \sum_{\underline{t} \in \calT_k} \prod_{h = 1}^{|\underline{t}|} \int_{\R_+} e^{-\gamma} \gamma^{n_h} \rho(\dd \gamma),
\end{aligned}
\end{equation*}
where the last equality is just due to a change in notation. Therefore, \eqref{eq:janossy_result_specialized} can be written as
\begin{equation*}
\begin{aligned}
 &  j_k(x_1,\ldots,x_k) =  k!\prob(N=k) \\
    &\quad \times \sum_{\underline{t} \in \calT_k} \left[ \prod_{h=1}^{|\underline{t}|} \int_{\X} \prod_{j: t_j = h} \kappa(x_i ; \theta) G_0(\dd \theta)\right] \frac{\exp\left\{ - \int_{\R_+} ( 1 - e^{-\gamma}) \rho(\dd \gamma) \right\} \prod_{h=1}^{|\underline{t}|}  \int_{\R_+} e^{-\gamma} \gamma^{n_h} \rho(\dd \gamma)}{k!\prob(N=k)},  
\end{aligned}
\end{equation*}
where the fraction term corresponds to $\prob(\underline{T}= \underline{t})$, where the law of $\underline{T}$ is defined in the statement. In conclusion, the thesis follows by expressing the last equation as an expected value with respect to $\underline{T}$.
\end{proof}

\end{document}
